\documentclass[12pt,a4paper]{article}

\usepackage[T1]{fontenc}
\usepackage[utf8]{inputenc}
\usepackage[romanian]{babel}
\usepackage{geometry}

\geometry{margin=2.5cm}

\title{Tutoriat 2\\Cautare Informata. Euristici. Greedy si A*}
\author{}
\date{}

\begin{document}

\maketitle

\section{Introducere}

În multe probleme de inteligență artificială, spațiul de stări este foarte mare, iar strategiile de căutare neinformată devin ineficiente.

Strategiile de căutare neinformată folosesc doar informațiile disponibile în definiția problemei și nu au cunoștințe despre cât de aproape se află un nod de soluție.

Pentru a explora mai eficient spațiul de căutare folosim căutarea informată, care utilizează cunoștințe suplimentare despre problemă pentru a ghida alegerea nodurilor explorate.

\section{Recapitulare. Probleme de căutare}

O problemă de căutare este definită prin:

\begin{itemize}
\item starea inițială;
\item operatorii (acțiunile disponibile);
\item testul scop;
\item funcția de cost.
\end{itemize}

O soluție reprezintă un drum de la starea inițială la o stare care satisface testul scop.

Exemplu:

Stare inițială: Arad

Scop: București

Acțiuni: deplasări între orașe conectate

Cost: distanța parcursă

Soluția optimă este drumul cu cost total minim.

\section{Frontiera și strategia de căutare}

Procesul de căutare construiește un arbore de căutare.

Frontiera reprezintă colecția de noduri care au fost generate, dar care nu au fost încă expandate.

Strategia de căutare specifică ordinea în care sunt explorate nodurile din frontieră.

Algoritmul general de căutare:

1. Inițializează frontiera cu starea inițială.

2. Cât timp frontiera nu este vidă:
   - selectează un nod conform strategiei;
   - dacă nodul este scop, întoarce soluția;
   - generează succesorii;
   - adaugă succesorii în frontieră.

\section{Evaluarea strategiilor de căutare}

O strategie de căutare este analizată după mai multe criterii.

Completitudine

Strategia găsește o soluție dacă soluția există.

Optimalitate

Strategia găsește cea mai bună soluție disponibilă.

Complexitate temporală

Timpul necesar pentru găsirea soluției.

Complexitate spațială

Cantitatea de memorie utilizată.

\section{Căutare neinformată}

Căutarea neinformată nu folosește nicio informație suplimentară despre poziția scopului.

Agentul poate:

\begin{itemize}
\item genera succesorii unei stări;
\item verifica dacă o stare este scop.
\end{itemize}

Exemple:

\begin{itemize}
\item Breadth First Search (BFS);
\item Depth First Search (DFS);
\item Uniform Cost Search (UCS);
\item Iterative Deepening Search (IDDFS).
\end{itemize}

Problema principală este că aceste strategii pot explora foarte multe stări inutile.

\section{Căutare informată}

Căutarea informată utilizează informații suplimentare despre problemă.

Ideea este să explorăm mai întâi stările care par mai apropiate de scop.

Aceste informații suplimentare sunt codificate cu ajutorul funcțiilor euristice.

\section{Euristici}

O euristică este o funcție:

h(n)

care estimează costul rămas de la nodul curent până la o stare scop.

O euristică nu calculează costul exact.

Ea oferă doar o estimare.

Proprietăți:

\begin{itemize}
\item $h(n) \ge 0$;
\item h(scop) = 0;
\item este specifică problemei;
\item oferă informații despre apropierea de soluție.
\end{itemize}

\section{Exemplu de euristică pentru harta României}

Pentru fiecare oraș putem defini:

h(x) = distanța în linie dreaptă dintre orașul x și București.

Exemple:

h(Arad) = 366

h(Sibiu) = 253

h(București) = 0

Această valoare nu reprezintă costul real al drumului, ci doar o estimare.

\section{Exemplu de euristică pentru 8-Puzzle}

Pentru problema 8-puzzle putem utiliza mai multe euristici.

Prima euristică:

h1 = numărul pieselor aflate pe poziții greșite.

A doua euristică:

h2 = suma distanțelor Manhattan ale tuturor pieselor până la pozițiile finale.

În practică, h2 oferă de obicei estimări mai bune decât h1.

\section{Problema clătitelor}

Problema sortării clătitelor constă în obținerea unei stive ordonate prin întoarcerea repetată a unor prefixe.

O posibilă euristică este:

h(x) = cea mai mare clătită care nu este încă pe poziția corectă.

Valoarea euristicii oferă o estimare a distanței până la configurația scop.

\section{Căutarea Greedy}

Ideea principală a algoritmului Greedy este:

explorează întotdeauna nodul care pare cel mai apropiat de scop.

Funcția de evaluare este:

f(n) = h(n)

Algoritmul ignoră costul parcurs până la nodul curent.

Contează doar estimarea până la scop.

\section{Exemplu intuitiv}

Presupunem că avem două orașe candidate.

Primul oraș are:

h = 150

Al doilea oraș are:

h = 80

Greedy îl va alege pe al doilea deoarece pare mai apropiat de București.

Nu contează cât a costat drumul până la acel oraș.

\section{Avantajele algoritmului Greedy}

\begin{itemize}
\item găsește rapid o soluție;
\item utilizează informația euristică;
\item explorează mai puține noduri decât BFS sau DFS.
\end{itemize}

\section{Dezavantajele algoritmului Greedy}

Nu garantează soluția optimă.

Exemplu:

Greedy poate găsi drumul:

Arad → Sibiu → Făgăraș → București

deși drumul optim este:

Arad → Sibiu → Râmnicu Vâlcea → Pitești → București

Prin urmare, Greedy poate produce soluții suboptimale.

\section{Proprietățile algoritmului Greedy}

Completitudine:

DA, dacă evităm ciclurile.

Optimalitate:

NU.

Complexitate timp:

$O(b^m)$

Complexitate spațiu:

$O(b^m)$

unde:

b = factorul de ramificare

m = adâncimea maximă explorată

\section{Uniform Cost Search}

Uniform Cost Search prioritizează nodurile cu cost acumulat minim.

Funcția utilizată este:

f(n) = g(n)

unde:

g(n) reprezintă costul real al drumului de la start la nodul curent.

UCS este optimal deoarece explorează nodurile în ordine crescătoare a costului.

\section{Interpretarea lui g(n)}

g(n) răspunde la întrebarea:

"Cât am plătit până acum pentru a ajunge aici?"

Această valoare este exactă și nu implică estimări.

\section{Interpretarea lui h(n)}

h(n) răspunde la întrebarea:

"Cât estimez că mai am până la scop?"

Această valoare este aproximativă.

\section{Algoritmul A*}

A* combină avantajele strategiilor UCS și Greedy.

Funcția de evaluare devine:

f(n) = g(n) + h(n)

unde:

g(n) = costul real parcurs

h(n) = costul estimat rămas

Algoritmul selectează nodul cu valoarea f minimă.

\section{Intuiția algoritmului A*}

Greedy folosește doar viitorul estimat.

UCS folosește doar trecutul.

A* utilizează simultan:

\begin{itemize}
\item trecutul exact;
\item viitorul estimat.
\end{itemize}

Prin această combinație, algoritmul este mai eficient decât UCS și mai sigur decât Greedy.

\section{Când se oprește algoritmul A*}

Nu ne oprim în momentul în care găsim pentru prima dată un nod scop.

Ne oprim atunci când nodul scop este extras din frontieră și nu mai există noduri care pot conduce la o soluție mai bună.

\section{Euristici admisibile}

O euristică este admisibilă dacă nu supraestimează niciodată costul real.

Formal:

$h(n) \le h^{*}(n)$

unde:

h*(n) reprezintă costul real minim până la scop.

Altfel spus:

estimarea trebuie să fie optimistă.

\section{Euristici inadmisibile}

O euristică este inadmisibilă dacă supraestimează costul real.

Consecință:

algoritmul A* poate pierde soluția optimă.

\section{Exemple de euristici admisibile}

Problema misionarilor și canibalilor:

h(n) = numărul persoanelor rămase pe malul inițial.

Problema 8-puzzle:

h(n) = numărul pieselor aflate pe poziții greșite.

h(n) = suma distanțelor Manhattan.

\section{Dominanța euristicilor}

Considerăm două euristici admisibile:

ha

și

hc

Dacă:

$ha(n) \ge hc(n)$

pentru toate stările,

atunci ha domină hc.

O euristică dominantă oferă de obicei performanțe mai bune deoarece este mai aproape de valoarea reala.

\section{Proprietățile algoritmului A*}

Completitudine:

DA

Optimalitate:

DA, dacă euristica este admisibilă.

Complexitate timp:

$O(b^m)$

Complexitate spațiu:

$O(b^s)$

unde:

b = factorul de ramificare

m = adâncimea maximă explorată

s = adâncimea soluției

\section{Comparație între UCS, Greedy și A*}

UCS

\begin{itemize}
\item optimal;
\item complet;
\item explorează multe noduri.
\end{itemize}

Greedy

\begin{itemize}
\item rapid;
\item folosește euristici;
\item nu este optimal.
\end{itemize}

A*

\begin{itemize}
\item folosește euristici;
\item este complet;
\item este optimal cu euristici admisibile;
\item explorează mai puține noduri decât UCS.
\end{itemize}

\section{Concluzii}

Căutarea informată utilizează informații suplimentare despre problemă pentru a ghida explorarea spațiului de stări.

Funcțiile euristice estimează distanța până la scop și permit reducerea numărului de noduri explorate.

Algoritmul Greedy explorează nodurile care par cele mai apropiate de scop și obține rapid soluții, dar nu garantează optimalitatea.

Algoritmul A* combină costul real și estimarea euristică prin funcția:

f(n) = g(n) + h(n)

și reprezintă una dintre cele mai importante tehnici de căutare informată, fiind complet și optimal atunci când utilizează euristici admisibile.

\end{document}